\begin{userbox}
Design a mix for this print job: 600 mm high, 1.5 m long wall, two filaments, 25 mm nozzle, 50 mm/s, open time 60 min. Space: OPC required 0.5-1.0, silica fume 0-0.15, limestone powder 0-0.3, w/b 0.28-0.42, s/b 0.8-1.8, PCE 0.1-1.2 \%, cellulose VMA 0-0.4 \%. The mix must also reach 30 MPa at 28 days. Run without asking.
\end{userbox}
\begin{toolbox}{describe\_vocabulary, 0.0 s}
\begin{lstlisting}[style=argstyle]
{"family": "all"}
\end{lstlisting}
\end{toolbox}
\begin{toolbox}{design\_for\_print\_job, 42.0 s}
\begin{lstlisting}[style=argstyle]
{"job": {"object_type": "wall", "target_height_mm": 600, "footprint_mm": 1500, "wall_filaments": 2, "nozzle_d_mm": 25, "print_speed_mm_s": 50, "open_time_min": 60}, "space": {"system_type": "mortar", "binder": {"portland_cement": {"lo": 0.5, "hi": 1, "required ...}
\end{lstlisting}
\end{toolbox}
\begin{toolbox}{retrieve\_literature, 0.6 s}
\begin{lstlisting}[style=argstyle]
{"system_type": "mortar", "prefer_3dcp": true, "n": 3, "targets": [{"quantity": "compressive_strength", "kind": "ge", "lo": 30, "unit": "MPa", "age_d": 28, "comparability_group": "comp_cube50"}]}
\end{lstlisting}
\end{toolbox}
\begin{toolbox}{predict\_properties, 1.0 s}
\begin{lstlisting}[style=argstyle]
{"mix": {"system_type": "mortar", "binder": {"portland_cement": 0.78, "limestone_powder": 0.22}, "w_b": 0.286, "s_b": 1.73, "admixtures_pct": {"superplasticiser_pce": 1.02}, "conditions": {"age_d": 28, "comparability_group": "comp_cube50", "curing": "moist", " ...}
\end{lstlisting}
\end{toolbox}
\begin{toolbox}{assess\_buildability, 0.6 s}
\begin{lstlisting}[style=argstyle]
{"job": {"object_type": "wall", "target_height_mm": 600, "footprint_mm": 1500, "wall_filaments": 2, "nozzle_d_mm": 25, "print_speed_mm_s": 50, "open_time_min": 60}, "mix": {"system_type": "mortar", "binder": {"portland_cement": 0.78, "limestone_powder": 0.22}, ...}
\end{lstlisting}
\end{toolbox}
\begin{agentbox}
I'll run the print-job design within your specified mix limits, then check buildability and the strength target. I'll report any defaults used for unspecified printing or test conditions.\par
The search found no candidate meeting the original screening criteria. Its best buildability candidate fails both plastic-collapse and free-wall buckling checks, and the required delay between layers exceeds the open-time allowance. I'm checking that candidate's strength and the material improvements needed before presenting it as a trial mix.\par
\textbf{No feasible mix was found for the specified job.} The search evaluated 4,096 mixes; candidates returned after automatic screening relaxation still failed physical buildability checks.\par
\par
The best-ranked buildability candidate, rounded and independently reassessed, is a \textbf{laboratory trial mix--not an approved printing mix}:\par
\par
\par\smallskip\begin{adjustbox}{max width=\linewidth}\begin{tabular}{p{0.47\linewidth}p{0.47\linewidth}}
Component or ratio & Trial value \\
OPC, binder mass fraction & 0.78 \\
Limestone powder, binder mass fraction & 0.22 \\
Silica fume & Omitted \\
w/b & 0.286 \\
s/b & 1.73 \\
PCE, product mass/binder mass & 1.02\% \\
\multicolumn{1}{l}{\textit{\ldots 1 further rows in the transcript}} \\
\end{tabular}\end{adjustbox}\par\smallskip
\par
Predictions are q50 [q10, q90]:\par
\par
\par\smallskip\begin{adjustbox}{max width=\linewidth}\begin{tabular}{p{0.47\linewidth}p{0.47\linewidth}}
Property & Prediction \\
Compressive strength, 28 days & \textbf{55.0 [34.2, 96.7] MPa} \\
Initial static yield stress & 563 [40, 9,654] Pa \\
Structuration rate & 0.778 [0.016, 14.605] Pa/s \\
Plastic viscosity & 6.73 [0.99, 48.03] Pa.s \\
\end{tabular}\end{adjustbox}\par\smallskip
\par
Strength is promising, but \textbf{plastic collapse governs}. The reassessment indicates collapse around layer 6 and free-wall buckling around layer 11, against 48 required layers. The minimum stability cycle is approximately \textbf{272 s}, while open time permits at most \textbf{75 s}: there is no feasible cycle-time window.\par
\par
Defaults used: layer height \textbf{12.5 mm}, filament width \textbf{29.61 mm}, \textbf{60 s} cycle, no dwell, printing immediately after mixing, density \textbf{2,100 kg/m\textasciicircum{}3}, safety factor \textbf{1.5}, and stiffness/yield-stress ratio \textbf{25}. Strength assumes moist curing and \texttt{comp\_cube50}.\par
\par
To progress, experimentally target \textbf{Athix >=3.52 Pa/s} at the predicted initial yield stress, or \textbf{initial static yield stress >=8,464 Pa} at the predicted structuration rate. These are alternative plastic-collapse requirements; buckling still needs bracing, returns, or greater wall thickness. VMA adjustment within your range is worth testing; an accelerator would require expanding the search space.\par
\par
[...]\par
\end{agentbox}